\section*{第 \ref{ch:b1:heat-budget} 章——地球热收支}

\begin{solution}{exo:b1:heat-budget:1}
$q = k\,\mathrm dT/\mathrm dz = 2.5\times\qty{0.022}{K.m^{-1}} =
\qty{0.055}{W.m^{-2}} = \qty{55}{mW.m^{-2}}$。通过 \qty{1}{km^2} $=
\qty{e6}{m^2}$：\qty{5.5e4}{W}，即 \qty{55}{kW}。
\end{solution}

\begin{solution}{exo:b1:heat-budget:2}
大陆 $\num{2.073e14}\times0.0709 = \qty{1.470e13}{W} = \qty{14.70}{TW}$；
海洋 $\num{3.028e14}\times0.1054 = \qty{31.92}{TW}$；总计 \qty{46.61}{TW}，并
平均\omterm{def:b1:heat-budget:heat-flow}{热流} $\num{4.661e13}/\num{5.101e14} = \qty{91.4}{mW.m^{-2}}$（
已发表的 46.7 和 91.6 之间的差异仅源于两个平均值的取整）。太阳输入的
占比为 $46.6/173000 = \num{2.7e-4}$，即
\qty{0.027}{\percent}。
\end{solution}

\begin{solution}{exo:b1:heat-budget:3}
铀：$0.99274\times\num{9.513e-5} + 0.00720\times\num{5.685e-4} =
\num{9.444e-5} + \num{4.09e-6} = \qty{9.85e-5}{W.kg^{-1}}$；其中稀有的
\ce{^{235}U} 贡献了其中的 \qty{4}{\percent}。对于钾，
\ce{^{40}K} 的质量分数为 $\num{1.17e-4}\times39.96/39.10 = \num{1.196e-4}$，因此
$\num{1.196e-4}\times\num{2.847e-5} = \qty{3.40e-9}{W.kg^{-1}}$。
\end{solution}

\begin{solution}{exo:b1:heat-budget:4}
$\qty{1}{\micro cal.cm^{-2}.s^{-1}} = \num{4.184e-6}\,\unit{J.s^{-1}}/\qty{e-4}{m^2}
= \qty{41.84}{mW.m^{-2}}$：美国东部 \qty{33.5}{mW.m^{-2}}，盆岭
区 \qty{58.6}{mW.m^{-2}}，内华达山脉 \qty{16.7}{mW.m^{-2}}。若 $A_0D$ 的单位为
\unit{mW.m^{-2}}（$A_0$ 以 \unit{\micro W.m^{-3}}、$D$ 以 \unit{km} 计）：
$33.5 + 3.0\times7.5 = \qty{56.0}{mW.m^{-2}}$，$58.6 + 3.0\times9.4 =
\qty{86.8}{mW.m^{-2}}$，$16.7 + 3.0\times10 = \qty{46.7}{mW.m^{-2}}$。
同一花岗岩在盆岭区给出的\omterm{def:b1:heat-budget:heat-flow}{热流}几乎是
内华达山脉的两倍：差异在于下方的\omterm{def:b1:layered-earth:layers}{地幔}。
\end{solution}

\begin{solution}{exo:b1:heat-budget:5}
$T(10\,\text{千米}) = 10 + \frac{0.060\times10^4}{2.5} - \frac{\num{2e-6}\times10^8}{2\times2.5}
= 10 + 240 - 40 = \qty{210}{\degreeCelsius}$。下方，$q_m = 60 - 2.0\times10 =
\qty{40}{mW.m^{-2}}$，梯度为 \qty{16}{K.km^{-1}}，并且
$T(35\,\text{千米}) = 210 + 16\times25 = \qty{610}{\degreeCelsius}$。
\end{solution}

\begin{solution}{exo:b1:heat-budget:6}
页岩：梯度 $12.0/0.300 = \qty{40}{K.km^{-1}}$，$q = 1.5\times0.040 =
\qty{60}{mW.m^{-2}}$。花岗岩：$10.0/0.500 = \qty{20}{K.km^{-1}}$，
$q = 3.0\times0.020 = \qty{60}{mW.m^{-2}}$。$R(300) = 300/1.5 =
\qty{200}{m^2.K.W^{-1}}$，$R(800) = 200 + 500/3.0 = \qty{366.7}{m^2.K.W^{-1}}$。
三点 $(0, 12.0)$、$(200, 24.0)$、$(366.7, 34.0)$ 共线，
斜率为 $22.0/366.7 = \qty{0.0600}{W.m^{-2}}$：\omterm{def:b1:heat-budget:heat-flow}{热流}相同，
两种岩石中均为 \qty{60}{mW.m^{-2}}；只有梯度改变，且与
传导率成反比。
\end{solution}

\begin{solution}{exo:b1:heat-budget:7}
$\tau = L^2/\kappa$：\qty{1}{m}，\qty{e6}{s} 或 \num{12}~天；\qty{10}{m}，
\qty{e8}{s} 或 \num{3.2}~年；\qty{100}{km}，\qty{e16}{s} 或
\qty{317}{Ma}；\qty{2900}{km}，\qty{8.4e18}{s} 或 \qty{267}{Ga}。在
\qty{10}{m} 处完全感觉不到日循环，而年循环的周期短于
穿过 \qty{10}{m} 所需的三年，其影响极其微弱：只有
持续数年或更久的变化才能传递过去。最后一个时间是
地球年龄的 $267/4.54 = 59$ 倍。
\end{solution}

\begin{solution}{exo:b1:heat-budget:8}
(a) $1646 - 0.54\times50 = \qty{1619}{K}$，即 \qty{1346}{\degreeCelsius}，
已发表的 \qty{1350}{\degreeCelsius} 为取整值。(b) 平均梯度
$(0.54 + 0.36)/2 = \qty{0.45}{K.km^{-1}}$，因此 $1646 + 0.45\times360 =
\qty{1808}{K}$，比 \qty{1799}{K} 高出 \qty{9}{K}（\qty{0.5}{\percent}）：
梯度并非线性下降，但估计仍然接近。
\end{solution}

\begin{solution}{exo:b1:heat-budget:9}
(a) 由 $G^2 = T_0^2/(\pi\kappa t)$ 可得 $t = T_0^2/(\pi\kappa G^2)$。
(b) $G = 0.02\times\tfrac59\,\unit{K}/\qty{0.3048}{m} = \qty{0.03645}{K.m^{-1}} =
\qty{36.45}{K.km^{-1}}$; $t = 3900^2/(\pi\times\num{1.2e-6}\times0.03645^2) =
\num{1.521e7}/\num{5.009e-9} = \qty{3.04e15}{s}$，即 \qty{96}{Ma}，
开尔文的结果。(c) $t = 1300^2/(\pi\times\num{e-6}\times0.036^2) =
\qty{4.15e14}{s} = \qty{13}{Ma}$：更为现实的初始温度使导热地球
更加年轻。
\end{solution}

\begin{solution}{exo:b1:heat-budget:10}
(a) 侵蚀之后，新地表位于原深度 $e$ 处，该处的
生热率为 $A_0' = A_0e^{-e/D}$，\omterm{def:b1:heat-budget:heat-flow}{热流}为 $q(e) = q_m +
A_0De^{-e/D} = q_m + A_0'D$；从新地表量起，生热率
仍按 $A_0'e^{-z'/D}$ 递减。深成岩体沿同一直线向下移动。
(b) 此时层厚为 $h - e$，生热率 $A$ 不变，因此 $q_0' = q_m + A(h - e)$：
即斜率为 $h - e$、过 $q_m$ 的直线。(c) 只有指数模型：在
均匀层模型中，经不同深度侵蚀的深成岩体将落在斜率不同的直线上，
同一省的数据点将会散乱。观测到单一直线，
这一事实支持生热率大致按指数规律随
深度递减。
\end{solution}

\begin{solution}{exo:b1:heat-budget:11}
(a) $\tau_P = RL/(3\kappa)$：对于 \qty{50}{km}，$\num{6.4e6}\times\num{5e4}/\num{3e-6}
= \qty{1.07e17}{s} = \qty{3.4}{Ga}$；对于 \qty{100}{km}，\qty{6.8}{Ga}。
(b) $q = kT_i/L = 3\times1300/\num{5e4} = \qty{78}{mW.m^{-2}}$，开始时
而 $78\,e^{-2/3.38} = \qty{43}{mW.m^{-2}}$，经过 \qty{2}{Ga} 后。(c)
$3\times1300/\sqrt{\pi\times\num{e-6}\times\num{6.31e16}} = 3900/\num{4.45e5} =
\qty{8.8}{mW.m^{-2}}$：充分混合的地球
经二十亿年仍多失去五倍的热量。其整个内部为地表供热，而在导热地球中只有
外层壳体如此。
\end{solution}

\begin{solution}{exo:b1:heat-budget:12}
(a) $\sqrt{\pi\kappa\times\qty{1}{Ma}} = \sqrt{\pi\times\num{e-6}\times\num{3.156e13}}
= \qty{9957}{m}$，因此 $C = 3\times1336/9957 = \qty{0.403}{W.m^{-2}}$，对应
$t$ 以 Ma 计，即 \qty{403}{mW.m^{-2}.Ma^{1/2}}，为拟合值
490 的 \qty{82}{\percent}。模型需要更大的 $k\,\Delta T$（约 \qty{4900}{W.m^{-1}}）：
\omterm{def:b1:layered-earth:layers}{地幔}传导率更高，或内部温度更高，即沿绝热线的情形。
(b) $490/\sqrt1 = 490$，$490/\sqrt{10} = 155$ 及 $490/\sqrt{100} =
\qty{49}{mW.m^{-2}}$。(c) $z = 2.326\sqrt{\num{e-6}\times\num{3.156e15}} =
2.326\times\qty{56180}{m} = \qty{131}{km}$。
\end{solution}

\begin{solution}{pb:b1:heat-budget:1}
\textbf{1.} 梯度 $= (T_{500} - T_{100})/\qty{0.4}{km}$：A \qty{15.25}{K.km^{-1}}，
B \num{19.50}，C \num{19.50}，D \num{25.50}，E \qty{28.25}{K.km^{-1}}。

\textbf{2.} 钻探会扰动岩石的温度（钻井液使
井壁降温或升温），这种扰动需要数周至数月才会消散。
顶部几十米的温度随季节循环及近期地表
温度变化而变化。

\textbf{3.} $q_0 = k\times$梯度：A $2.9\times15.25 = \qty{44.225}{mW.m^{-2}}$，
B \qty{50.70}{mW.m^{-2}}，C \qty{60.45}{mW.m^{-2}}，D \qty{68.85}{mW.m^{-2}}，
E \qty{84.75}{mW.m^{-2}}。

\textbf{4.} A 至 D 均低于大陆平均值，E 高出
\qty{13.9}{mW.m^{-2}}，即 \qty{20}{\percent}。

\textbf{5.} \omterm{def:b1:heat-budget:heat-flow}{热流}等于传导率乘以梯度：C 的传导率更高（3.1
对 2.6），因此同样的梯度携带的热量多 \qty{19}{\percent}。

\textbf{6.} 等温线只能部分地随地形起伏：山顶之下，
它们的抬升幅度小于地表，因此在其下方彼此间距更大，
实测梯度偏低（山谷之下则偏高）。
比较深成岩体之前，需要先对地形作校正。

\textbf{7.} $\rho C_U H_U$：A $2650\times\num{2.0e-6}\times\num{9.85e-5} =
\qty{0.522}{\micro W.m^{-3}}$，B 0.914，C 1.305，D 1.827，E 2.610。

\textbf{8.} 钍：A 0.558，B 0.976，C 1.255，D 1.742，E 2.439。

\textbf{9.} 钾：A $2650\times0.030\times\num{3.40e-9} = 0.270$，B 0.315，
C 0.360，D 0.378，E 0.405。

\textbf{10.} $A_0$ = 1.35、2.20、2.92、3.95 及 \qty{5.46}{\micro W.m^{-3}}。
钾的贡献最小，从 A 中占总量的 \qty{20}{\percent} 降至
E 中的 \qty{7}{\percent}；铀和钍几乎平分其余部分。

\textbf{11.} 是：五个数据点稳步上升，与一条直线的偏差约在
\qty{1.5}{mW.m^{-2}} 以内。

\textbf{12.} $\bar A_0 = 15.88/5 = 3.176$，$\bar q_0 = 308.975/5 = 61.795$；
$\sum(A_0 - \bar A_0)^2 = 10.1681$；$\sum(A_0 - \bar A_0)(q_0 - \bar q_0) =
101.146$。

\textbf{13.} $D = 101.146/10.1681 = 9.95$（单位为 \unit{mW.m^{-2}} 每
\unit{\micro W.m^{-3}}，即以千米计）：$D = \qty{9.95}{km}$，且 $q_m = 61.795 -
9.95\times3.176 = \qty{30.20}{mW.m^{-2}}$。

\textbf{14.} 斜率约为 \qty{10}{km}，即三个省份的斜率；
约化\omterm{def:b1:heat-budget:heat-flow}{热流} \qty{30}{mW.m^{-2}}，接近美国东部的
\qty{33}{mW.m^{-2}}：这是一个平静的省份，下方\omterm{def:b1:layered-earth:layers}{地幔}较冷，不同于
盆岭省。

\textbf{15.} $q_m/k = \qty{12}{K.km^{-1}}$ 及 $A_0D^2/k = \num{5.46e-6}\times\num{e8}/2.5
= \qty{218.4}{K}$：$T(z) = 10 + 12z + 218.4\,(1 - e^{-z/10})$，$z$ 以千米计，$T$ 以
\unit{\degreeCelsius} 计。

\textbf{16.} $q(z) = 30 + 54.6\,e^{-z/10}$：在 \qty{10}{km} 处为 \qty{50.1}{mW.m^{-2}}，
在 \qty{20}{km} 处为 \num{37.4}，在 \qty{35}{km} 处为 \qty{31.7}{mW.m^{-2}}。

\textbf{17.} $30/84.6 = \qty{35}{\percent}$：E 上方的\omterm{def:b1:heat-budget:heat-flow}{热流}中近三分之二
由其放射性产热层产生。

\textbf{18.} $T(10) = 10 + 120 + 218.4\times0.632 = \qty{268}{\degreeCelsius}$。

\textbf{19.} $T(20) = 10 + 240 + 218.4\times0.865 = \qty{439}{\degreeCelsius}$。

\textbf{20.} 在 A 之下，$A_0D^2/k = 1.35\times100/2.5 = \qty{54.0}{K}$：
$T(35) = 10 + 420 + 54.0\times0.970 = \qty{482}{\degreeCelsius}$.

\textbf{21.} 由 $A_0D^2/k$ 可得：A 之下为 \qty{54}{K}，E 之下为 \qty{218}{K}。深部
\omterm{def:b1:heat-budget:heat-flow}{热流}相同；E 的\omterm{def:b1:layered-earth:layers}{地壳}中热源多出四倍，这些热源产生的热量
必须沿同一岩层向上传导，这就需要
整段更陡的梯度。

\textbf{22.} $1300 = 10 + 218.4 + 12z$ 给出 E 之下 $z = \qty{89.3}{km}$；
$1300 = 10 + 54.0 + 12z$ 给出 A 之下 \qty{103.0}{km}。

\textbf{23.} 两者均落在大陆\omterm{def:b1:layered-earth:lithosphere}{岩石圈}的范围之内，
即年轻地区之下约一百千米或更浅，至克拉通之下约 \qty{250}{km}，
靠近其较薄的一侧：在约化\omterm{def:b1:heat-budget:heat-flow}{热流}为 \qty{30}{mW.m^{-2}}、这一传导率的条件下，
传导型\omterm{def:b1:heat-budget:geotherm}{地温线}约在
\qty{100}{km} 处达到\omterm{def:b1:layered-earth:layers}{地幔}温度；在放射性更强的深成岩体之下，则早 \qty{14}{km} 达到。

\textbf{24.} $T(35) = 10 + 420 + 218.4\times(1 - e^{-3.5}) = 430 + 211.8$：深成岩体 E
之下的\omterm{def:b1:layered-earth:discontinuity}{莫霍面}位于
\textbf{\qty[text-series-to-math=true]{642}{\degreeCelsius}}.
\end{solution}
